\section{Experiment: Effectiveness and Learning Analysis}
In our first experiment, we evaluate BanditGPT's ability to identify and converge on optimal models in a high-quality routing scenario. This experiment serves as the primary benchmark for learning efficiency and competitive performance against established baselines.

\subsection{Objective Alignment (Profile Configuration)}
For the primary effectiveness comparison (Regret), we configure BanditGPT with a \textbf{Quality-Dominant Profile} ($w_{quality}=0.97, w_{cost}=0.02, w_{latency}=0.01$). This decision is grounded in control theory:
\begin{enumerate}
    \item \textbf{Metric Consistency:} The Oracle reward $r^*$ is derived from LMSYS preference labels, which measure pure response quality. To minimize Regret ($R_T = \sum r^* - r_t$), the agent's utility function must align with this ground truth.
    \item \textbf{Isolating Learning Efficiency:} Using a cost-sensitive profile (e.g., ``Arbitrage'') would introduce a strategic bias---intentionally sacrificing reward $r_t$ for cost savings. While practically useful, this confounds the measurement of learning efficiency, as it becomes impossible to distinguish between a ``Smart Trade-off'' and a ``Routing Error.''
\end{enumerate}
By targeting Max Quality, we rigorously test the bandit's ability to identify and converge on the optimal model solely through feature-based generalization.

\subsection{Impact of Prior Calibration}
Figure \ref{fig:heatmap} visualizes the hyperparameter landscape derived from our 3-Fold Cross-Validation on the development set. We observe two distinct regimes:
\begin{enumerate}
    \item \textbf{Exploration Suppression:} Performance degrades consistently as the exploration rate $\alpha$ increases (left-to-right gradient). This validates our hypothesis that high-quality procedural priors effectively "pre-solve" the exploration problem, rendering high-variance sampling strategies (like $\epsilon$-greedy or high-$\alpha$ LinUCB) unnecessary and costly.
    \item \textbf{Plasticity Sweet Spot:} We identify a convex basin of optimality at $N_{eff}=5.0$. While the router is robust to stiffer priors ($N_{eff}=100$ achieves competitive regret of 0.0158), the global minimum at $N_{eff}=5.0$ (Regret 0.0150) suggests that treating priors as a "strong nudge" rather than "dogma" allows the agent to correct minor Sim-to-Real mismatches most efficiently.
\end{enumerate}
Based on this analysis, we fix $N_{eff}=5.0$ and $\alpha=0.05$ for all subsequent evaluations on the held-out test set.

\subsection{Regret Analysis and Baseline Comparison}
\begin{figure}[t]
    \centering
    \includegraphics[width=0.9\linewidth]{experiments/01_effectiveness/results/effectiveness_plot.png}
    \caption{\textbf{Cumulative Regret (10 Seeds).} Mean cumulative regret across 10 independent trials on the held-out test set ($n=800$). Shading represents $\pm 1$ standard deviation. BanditGPT exhibits a steep learning curve in the first 100 steps followed by stable convergence, outperforming both static and adaptive baselines.}
    \label{fig:regret_comparison}
\end{figure}

\begin{table}[h]
\centering
\caption{\textbf{Effectiveness Comparison Results.} Cumulative regret (Mean $\pm$ 95\% CI) across 10 independent trials. BanditGPT significantly outperforms all baselines, achieving a $\sim$3$\times$ reduction in regret compared to the strongest vanilla baseline.}
\label{tab:effectiveness_results}
\begin{tabular}{llr}
\hline
\textbf{Category} & \textbf{Method} & \textbf{Cumulative Regret} \\ \hline
Non-Learning & Random & $29.34 \pm 11.18$ \\
SOTA Baseline & RouteLLM (Matrix Factorization) & $24.03 \pm 8.13$ \\
Ablation & BanditGPT (HLE-Only) & $23.78 \pm 9.02$ \\
Standard Bandit & Vanilla LinUCB & $20.63 \pm 5.69$ \\
\textbf{Proposed} & \textbf{BanditGPT (Procedural Warmup)} & $\mathbf{6.95 \pm 3.18}$ \\ \hline
\end{tabular}
\end{table}

We compare BanditGPT against three baselines using Oracle Rewards:
\begin{itemize}
    \item \textbf{RouteLLM (SOTA):} The Matrix Factorization router from \cite{ong2024routellm}, using the official library with mapped model IDs.
    \item \textbf{Vanilla LinUCB:} A standard disjoint LinUCB policy ($d=1$, bias only). We set $\alpha=1.0$ to ensure sufficient exploration in the absence of informative priors.
    \item \textbf{BanditGPT (HLE):} Our model initialized only with diagonal priors (HLE scores) but no covariance structure ($N_{eff}=10$), serving as an ablation for the Procedural Warmup.
\end{itemize}

The results in Table \ref{tab:effectiveness_results} demonstrate that BanditGPT's procedural warmup provides a critical performance leap. While standard baselines struggle with the cold-start problem, our model identifies optimal routing strategies significantly faster, resulting in the lowest terminal regret and highest stability (as indicated by the narrowest confidence intervals among learning methods).
